
\subsection{Model Architecture}\label{sec:architecture}

We use an encoder-decoder architecture for sequence to sequence learning \citep{sutskever2014sequence} and we closely follow the architectural choices presented in \citet{vaswani2017transformer}.
Our self-attention baseline is the fairseq re-implementation of the Transformer Big architecture  \citep{ott2018scaling}.\footnote{\url{https://github.com/pytorch/fairseq}}

The encoder and decoder networks have $N$ blocks each.
Encoder blocks contain two sub-blocks: The first is a self-attention module (\textsection\ref{sec:background}), a \sdconv{} module (\ref{sec:sdconv}), or a \tvsdconv{} module (\textsection\ref{sec:tvsdconv}).
The second sub-block is a feed-forward module: $ReLU(\emW^1 \emX + b_1) \emW^2 + b_2$ where $\emW^1 \in \R^{d \times d_{ff}}$, $\emW^2 \in \R^{d_{ff} \times d}$ and $d=1024$, $d_{ff}=4096$ unless otherwise stated.
Sub-blocks are surrounded by residual connections \citep{he2015deep} and layer normalization \citep{ba2016layer}. 

Decoder blocks are identical except that they have an additional source-target attention sub-block between the self-attention and feed-forward module. 
The source-target attention is equivalent to the self-attention module, except that the values and keys are projections over the encoder output for each source word.

Words are fed to the encoder and decoder networks in $d$ dimensional embeddings. We add sinusoidal position embeddings to encode the absolute position of each word in the sequence \citep{kaiser2017depthwise,vaswani2017transformer}.
The model computes a distribution over vocabulary $V$ by transforming the decoder output via a linear layer with weights $\emW^V \in \R^{d \times V}$ followed by softmax normalization. 

\sdconv{} and \tvsdconv{} are identical to Transformer Big, except that self-attention modules are swapped with either fixed or dynamic convolutions.
These models also use fewer parameters per block (cf. Figure~\ref{fig:sdconv} and Figure~\ref{fig:tvsdconv}) and we 
therefore increase the number of blocks to $N=7$ for the encoder to roughly match the parameter count of Transformer Big.
We generally set $H=16$.
Both \sdconv{} and \tvsdconv{} set the the encoder and decoder kernel sizes to 3, 7, 15, 31x4 for each block respectively; except for the decoder where we have only three top layers with kernel size 31.




\subsection{Datasets and Evaluation}\label{sec:datasets}

To get a thorough understanding of the limitations of \sdconv{} and \tvsdconv{} we evaluate on three different tasks: machine translation, language modeling and abstractive summarization.

\noindent {\bf Machine Translation.}
We report results on four benchmarks: 
For WMT English to German (En-De) we replicate the setup of~\citet{vaswani2017transformer}, based on WMT'16 training data with 4.5M sentence pairs, we validate on newstest2013 and test on newstest2014.\footnote{\scriptsize \url{https://github.com/tensorflow/tensor2tensor/blob/321bacaa3abcca5dbf341ed6fb3d4a1531e513ff/tensor2tensor/data_generators/translate_ende.py\#L60-L63}}
The vocabulary is a 32K joint source and target byte pair encoding (BPE; \citealt{sennrich2016bpe}).
For WMT English to French (En-Fr), we borrow the setup of~\citet{gehring2017convs2s} with 36M training sentence pairs from WMT'14, validate on newstest2012+2013 and test on newstest2014.
The 40K vocabulary is based on a joint source and target BPE factorization.

For WMT English to Chinese (Zh-En), we pre-process the WMT'17 training data following \citet{hassan2018parity} resulting in 20M sentence pairs. 
We develop on devtest2017 and test on newstest2017. 
For IWSLT'14 German-English (De-En) we replicate the setup of \citet{edunov2018classical} for 160K training sentence pairs and 10K joint BPE vocabulary.
For this benchmark only, data is lowercased.

For \ende{}, \enfr{}, we measure case-sensitive tokenized BLEU.\footnote{\scriptsize{\url{https://github.com/moses-smt/mosesdecoder/blob/master/scripts/generic/multi-bleu.perl}}}
For \ende{} only we apply compound splitting similar to~\citet{vaswani2017transformer}.
For \zhen{} we measure detokenized BLEU to be comparable to \citet{hassan2018parity}.\footnote{SacreBLEU hash: \texttt{\scriptsize
BLEU+case.mixed+lang.zh-en+numrefs.1+smooth.exp+test.wmt17+tok.13a+version.1.2.11}}

We train three random initializations of a each configuration and report test accuracy of the seed which resulted in the highest validation BLEU. 
Ablations are conducted on the validation set and we report the mean BLEU and standard deviation on this set.
\ende{}, \enfr{} are based on beam search with a beam width of 5, \iwslt{} uses beam 4, and \zhen{} beam 8 following \citet{hassan2018parity}.
For all datasets, we tune a length penalty as well as the number of checkpoints to average on the validation set.

\noindent {\bf Language Modeling.} 
We evaluate on the large-scale \gbw{} dataset \citep{chelba2013one} which contains 768M tokens and has a vocabulary of nearly 800K types. 
Sentences in this dataset are shuffled and we batch sentences independently of each other.
Models are evaluated in terms of perplexity on the valid and test portions.

\noindent {\bf Summarization.}
We test the model's ability to process long documents on the \cnndm{} summarization task \citep{hermann15,nallapati16} comprising over 280K news articles paired with multi-sentence summaries. 
Articles are truncated to 400 tokens \citep{see17} and we use a BPE vocabulary of 30K types \citep{fan17}. 
We evaluate in terms of F1-Rouge, that is Rouge-1, Rouge-2 and Rouge-L \citep{lin04}.\footnote{We use the following parameters for \texttt{ROUGE-1.5.5.pl}: -m -a -n 2}
When generating summaries, we follow standard practice in tuning the maximum output length, disallowing repeating the same trigram, and we apply a stepwise length penalty \citep{paulus17intra,fan17,wu2016google}. 



\subsection{Training and Hyperparameters}


\noindent {\bf Translation.}
We use a dropout rate of 0.3 for WMT En-De and IWSLT De-En, 0.1 for WMT En-Fr, and 0.25 for WMT Zh-En. 
WMT models are optimized with Adam and a cosine learning rate schedule \citep{kingma2015adam,loshchilov2016cosine} where the learning rate is first linearly warmed up for 10K steps from $10^{-7}$ to $10^{-3}$ and then annealed following a cosine rate with a single cycle. 
For IWSLT'14 De-En, we use a schedule based on the inverse square root of the current step \citep{vaswani2017transformer}.
We train the WMT models on 8 NVIDIA V100 GPUs for a total of 30K steps on WMT En-De, 40K steps for WMT Zh-En and 80K steps for WMT En-Fr. 
For IWSLT De-En we train for 50K steps on a single GPU.

We use floating point 16 precision and accumulate the gradients for 16 batches before applying an update \citep{ott2018scaling}, except for IWSLT where we do not accumulate gradients.
Batches contain up to 459K source tokens and the same number of target tokens for both WMT En-De and WMT Zh-En, 655K for En-Fr, and 4K for IWSLT De-En.
We use label smoothing with $0.1$ weight for the uniform prior distribution over the vocabulary  \citep{szegedy2015inception,pereyra2017regularize}.

\noindent {\bf Language Modeling.} 
We follow the same setup as for translation but remove the encoder module.
For the \gbw{} benchmark we use an adaptive softmax output layer to reduce the computational burden of the large vocabulary \citep{grave2016arxiv,press2017using} and tie it with variable sized input word embeddings (Anonymous et al., 2018). 
The first 60K types in the adaptive softmax have dimension $1024$, the 100K types dimension $256$, and the last 633K types have size $64$.

We train on 32 GPUs with batches of 65K tokens for 975K updates. As optimizer we use Nesterov's accelerated gradient method \citep{sutskever2013icml} with a momentum value of $0.99$ and we re-normalize gradients if their norm exceeds $0.1$ \citep{pascanu2013difficulty}.
The learning rate is linearly warmed up from $10^{-7}$ to $1$ for 16K steps and then annealed using a cosine learning rate schedule \citep{loshchilov2016cosine} with one cycle.

\noindent {\bf Summarization.}
We train with Adam using the cosine learning rate schedule with a warmup of 10K steps and a period of 20K updates. 
We use weight decay 1e-3 and dropout 0.3. %, and train with label smoothing 0.1. 
